\section{Eksperimentalna analiza: plan i metodologija}
\label{sec:eksperimenti}
Eksperimentalni deo obuhvata istorijsko CPU poređenje i zaseban eksperiment
sa tri putanje AES-a na istoj platformi. Tekst definiše metodologiju;
rezultati još nisu prikupljeni za rad.
% TODO: Pre merenja usaglasiti protokol; tehničke smoke provere nisu rezultati članka.

\subsection{Pitanja i obuhvat eksperimenata}
Osnovni benchmark poredi DES, troključni TDEA i AES-128/192/256.
Novi eksperiment poredi softverski CPU AES-128, CPU AES-128 sa potvrđenim
AES-NI i NVIDIA GPU/CUDA putanju. Ispituje se zavisnost propusnosti i
latencije od veličine ulaza, bez pretpostavke da GPU mora biti brži.
% TODO: AES-192/256 na GPU-u dodati tek posle stabilne AES-128 implementacije i zasebne validacije.
% TODO: Definisati hipoteze: prednost CPU-a za male poruke, eventualni prag prednosti GPU-a i uticaj rezidentnih podataka. Dopustiti ishod bez preseka krivih.

\subsection{Okruženje i reproduktivnost}
Za svako izvršavanje beleže se CPU model, fizička i logička jezgra,
potvrđena AES putanja, GPU model, compute capability, slobodna i ukupna GPU
memorija, verzije NVIDIA drajvera, CUDA runtime-a, CuPy-ja, Pythona,
operativnog sistema i CPU kriptografske biblioteke. Uz seed, argumente,
identifikator verzije koda i hash CUDA izvora čuvaju se beleške o opterećenju.
% TODO: Zabeležiti napajanje, temperaturu, CPU/GPU taktove, WDDM/TCC režim i pozadinsko opterećenje.
% TODO: Razlikovati CUDA podršku drajvera prijavljenu kroz nvidia-smi od instaliranog Toolkit-a, NVRTC-a i stvarno učitanog runtime-a. Nepoznate vrednosti ostaviti nepoznatim.
% TODO: Seedovani ključevi i ponovljeni test ulazi služe isključivo sintetičkom eksperimentu.

\subsection{Osnovni CPU benchmark DES-a, TDEA i AES-a}
Postojeći CBC benchmark ostaje istorijsko poređenje bibliotečkih
implementacija. Vreme transformacije odvaja se od inicijalizacije šifre;
povratna dekripcija proverava se izvan merenog intervala. DES i TDEA ovde
nisu preporuke za novu zaštitu podataka.
% TODO: Zadržati warmup, više ponavljanja, sredinu, medijanu i SD; ulazi 1 KiB, 1 MiB i 10 MiB, opciono 100 MiB.
% TODO: Obeležiti da postojeći CPU benchmark koristi decimalne MB/s. Ne mešati njegove CBC rezultate sa novim ECB/CTR rezultatima u MiB/s.

\subsection{AES sa i bez CPU hardverske akceleracije}
\label{sec:aes-akceleracija}
CPU AES-NI koristi instrukcije namenjene AES transformacijama, dok
softverska putanja koristi prenosivu kompajliranu implementaciju iste
biblioteke~\cite{intel-aes-ni}. Eksperiment zahteva potvrđen izbor putanje;
sam zahtev \texttt{use\_aesni=True} nije dovoljan. Postojeće CBC poređenje
zadržava se, a za CPU--GPU eksperiment obe CPU putanje mere se ponovo
u istim ECB i CTR režimima kao GPU.
% TODO: Sačuvati proveru native inicijalizacije za konkretan mod i verziju biblioteke.
% TODO: Bez pouzdane kontrole AES-NI eksperiment označiti nepotpunim, ne preimenovati auto putanju u cpu_software.
% TODO: ARM akceleracija zahteva novi potvrđeni adapter; sadašnje oznake cpu_aesni odnose se na x86.
% TODO: Softverski AES na današnjem CPU-u nije rekonstrukcija performansi iz 2003.

\subsection{AES na NVIDIA GPU/CUDA platformi}
\label{sec:aes-gpu}
CUDA pristup raspoređuje obradu nezavisnih blokova na paralelne niti u SIMT
modelu~\cite{nvidia-cuda-model}. Python okruženje može upravljati korisničkim
kernelima~\cite{nvidia-cuda-python}; predviđena implementacija koristi
CuPy \texttt{RawKernel}, CUDA C izvor i eksplicitnu grid/block konfiguraciju
~\cite{cupy-rawkernel}. GPU paralelizacija i namenske CPU AES instrukcije
zato predstavljaju različite mehanizme ubrzanja.

ECB će služiti isključivo merenju nezavisne AES primitive, ne kao preporuka
za zaštitu podataka. Glavni paralelizabilni mod biće CTR:
\begin{equation}
C_i = P_i \oplus E_K(T_i).
\label{eq:eksperiment-ctr}
\end{equation}
Blokovi brojača mogu se obrađivati nezavisno. Nasuprot tome, CBC enkripcija
\begin{equation}
C_i = E_K(P_i \oplus C_{i-1})
\end{equation}
ima zavisnost između susednih blokova i nije primarni GPU encryption
benchmark~\cite{sp800-38a}. CTR sam ne obezbeđuje autentifikaciju.

Pre vremenskih merenja GPU rezultat mora odgovarati NIST SP 800-38A
vektorima F.1.1 i F.5.1 i referentnoj CPU implementaciji za isti ključ,
poruku i početni brojač. Proveravaju se jedan i više blokova, prenos
brojača preko bajtova i nepotpuni završni CTR blok. Neuspešna validacija
blokira benchmark; prolazak nije dokaz kriptografske bezbednosti.
% TODO: Implementirati i pregledati aes_gpu_kernel.cu. Skeleton trenutno namerno blokira kompilaciju; GPU ispravnost nije potvrđena.
% TODO: Dokumentovati poreklo CUDA implementacije; ne preuzimati nepoznatu PyPI šifru.
% TODO: Najpre čitljiva AES-128 implementacija, zatim validacija; optimizacije tek posle toga.

\subsection{CPU--GPU poređenje i uticaj veličine ulaza}
\label{sec:cpu-gpu-poredjenje}
Planirani ulazi su 1 KiB, 16 KiB, 1 MiB, 16 MiB, 100 MiB i 1 GiB.
Nedovoljna RAM/GPU memorija dovodi do preskakanja veličine uz zabeležen
razlog, ne do upisivanja nulte propusnosti. Za svaku kombinaciju moda,
veličine i putanje predviđeni su warmup i najmanje deset merenih ponavljanja;
broj nezavisnih sesija treba utvrditi pre studije.

Prvo CUDA izvršavanje može uključivati inicijalizaciju i kompilaciju
kernela; \texttt{RawKernel} može keširati prevedeni kod
~\cite{cupy-rawkernel}. JIT kompilacija, alokacija bafera i GPU priprema
ključa izuzimaju se iz stabilnog merenja. Vreme kernela meri se CUDA
događajima, uz obaveznu sinhronizaciju~\cite{cupy-performance}.
Posebno se beleže:
\begin{itemize}
\item \texttt{kernel\_ms}: GPU kernel između događaja;
\item \texttt{transfer\_ms}: zbir H2D i D2H zidnog vremena, uključujući
      sinhronizaciju i eventualno host staging kopiranje;
\item \texttt{end\_to\_end\_ms}: H2D, kernel, D2H i završna sinhronizacija
      za unapred alocirane bafere;
\item rezidentni prolaz: kernel i ukupno vreme poziva kada je ulaz već
      u GPU memoriji, bez transfera.
\end{itemize}
CPU redovi zasebno sadrže vreme kriptografske operacije i ukupno vreme
inicijalizacije šifre i enkripcije sa alokacijom izlaza.
Ovo je eksplicitno ograničen stabilni protokol, ne cold-start aplikativna
latencija. Razlike u alokacijama i pripremi ključa treba uzeti u obzir.
% TODO: Kao dopunski eksperiment kasnije meriti cold-start, alokacije i GPU key setup; ne pripisivati ove troškove sadašnjem end-to-end rezultatu.

Za vreme i throughput u MiB/s izračunavaju se sredina, medijana, uzoračka
standardna devijacija, minimum i maksimum. Faktori ubrzanja definišu se
odvojeno za GPU kernel i ukupnu GPU putanju:
\begin{equation}
S_{\mathrm{GPU/CPU}} =
\frac{T_{\mathrm{CPU}}}{T_{\mathrm{GPU}}}.
\end{equation}
U centralnoj tabeli koristiće se odnos medijana CPU AES-NI ukupnog vremena
i GPU ukupnog vremena. Vrednost veća od jedan označava kraće GPU vreme.
% TODO: Pronaći eventualni prag prednosti zasebno za ECB/CTR i kernel/ukupno vreme, uz varijabilnost između sesija. Ne interpolirati neizmereni prag kao činjenicu.
% TODO: Mešati redosled backend-a; za male poruke analizirati Python poziv, launch, tajmer, keš i sinhronizaciju. SD nije interval poverenja.
% TODO: Zameniti podacima za jedan jasno označen mod i run ID; -- znači nedostajuće.
\begin{table}[htbp]
\centering\small
\begin{tabularx}{\linewidth}{@{}l*{5}{>{\raggedright\arraybackslash}X}@{}}
\toprule
Veličina & CPU software & CPU AES-NI & GPU kernel & GPU end-to-end &
GPU/CPU AES-NI odnos \\
\midrule
1 KiB   & -- & -- & -- & -- & -- \\
16 KiB  & -- & -- & -- & -- & -- \\
1 MiB   & -- & -- & -- & -- & -- \\
16 MiB  & -- & -- & -- & -- & -- \\
100 MiB & -- & -- & -- & -- & -- \\
1 GiB   & -- & -- & -- & -- & -- \\
\bottomrule
\end{tabularx}
\caption{Obrazac CPU--GPU poređenja. Vremena su medijane u ms; CPU kolone
su ukupna vremena. Poslednja kolona je $T_{\mathrm{CPU\ AES-NI}}/
T_{\mathrm{GPU\ ukupno}}$. ECB i CTR popunjavati odvojeno.}
\label{tab:aes-cpu-gpu}
\end{table}


\subsection{Edukativna brute-force demonstracija}
% TODO: Zadržati ograničeni AES keyspace; ne pokretati puni DES/AES brute-force.
% TODO: Ekstrapolacije 2^56, 2^112, 2^128, 2^192 i 2^256 prikazati u naučnoj notaciji kao ilustraciju, uz razliku celog i očekivanog pretraživanja.
% TODO: GPU throughput enkripcije sa jednim ključem nije brzina pretraživanja različitih ključeva; ne koristiti ga za takvu procenu bez zasebnog eksperimenta. TDEA i MITM zahtevaju poseban model.

\subsection{Avalanche efekat}
% TODO: Menjati jedan bit poruke ili efektivnog ključa, isključiti DES paritet, dokumentovati TDEA degeneracije.
% TODO: Hamming distance normalizovati prema dužini bloka; prikazati raspodele, ne predstavljati test kao dokaz sigurnosti.

\subsection{Izvoz rezultata, grafikoni i ograničenja}
Sirovi CSV, statistički sažetak i JSON okruženja povezani su istim run ID-jem.
CPU polja za GPU kernel i transfer ostaju prazna, a ne nula.
Grafikoni i tabele generisaće se odvojeno za ECB i CTR, isključivo iz
validiranih izvršavanja. Ostatak ukupnog vremena nakon oduzimanja transfera
i kernela predstavlja procenu ostalog overhead-a, ne izolovano vreme
launch-a; različiti časovnici i rezolucija mogu dati mali negativan ostatak.

GPU kod je istraživački. Eksperiment ne ispituje constant-time svojstva,
bočne kanale, fault injection niti bezbedno upravljanje i brisanje ključeva
u GPU memoriji. Performanse ne dokazuju produkcionu bezbednost.
% TODO: Sačuvati verzije i sirova merenja; grafikone uključiti tek nakon pregleda, bez automatske zamene placeholdera.
\begin{figure}[htbp]
\centering
\fbox{\parbox[c][14mm][c]{0.90\linewidth}{\centering
Plan: \texttt{throughput\_vs\_size.pdf}\\
CPU software, CPU AES-NI, GPU kernel i GPU ukupno; logaritamska X osa.}}
\caption{Mesto za propusnost u zavisnosti od ulaza; rezultati još nisu uneti.}
\label{fig:gpu-throughput}
\end{figure}
\begin{figure}[htbp]
\centering
\fbox{\parbox[c][14mm][c]{0.90\linewidth}{\centering
Plan: \texttt{latency\_vs\_size.pdf}\\
Ukupna latencija tri putanje.}}
\caption{Mesto za poređenje latencije; rezultati još nisu uneti.}
\label{fig:gpu-latency}
\end{figure}
\begin{figure}[htbp]
\centering
\fbox{\parbox[c][14mm][c]{0.90\linewidth}{\centering
Plan: \texttt{gpu\_overhead\_breakdown.pdf}\\
Transfer, kernel i ostatak ukupnog vremena.}}
\caption{Mesto za strukturu GPU vremena; rezultati još nisu uneti.}
\label{fig:gpu-overhead}
\end{figure}
\begin{figure}[htbp]
\centering
\fbox{\parbox[c][14mm][c]{0.90\linewidth}{\centering
Plan: \texttt{speedup\_vs\_size.pdf}\\
Odvojeni kernel-only i end-to-end odnosi prema CPU AES-NI.}}
\caption{Mesto za faktore ubrzanja; ne pretpostavlja se prednost GPU-a.}
\label{fig:gpu-speedup}
\end{figure}

\subsection{Rezultati i diskusija}
% TODO: Pisati tek nakon validacije i stvarnih merenja. Razlikovati funkcionalnu proveru, pilot i rezultate za rad.
% TODO: Razmotriti ograničenja jednog računara, driver/WDDM overhead, prenos preko PCIe, efekte keša i memorijske kapacitete; ne tvrditi unapred da GPU pobeđuje.
